\chapter{고찰}

본 논문에서는 인간이 의사결정을 내릴 때 시행 횟수-이익 확률 모델 $\mathcal{M}_n(X, C)$를 고려한다는 가설을 제시하였다. 이 모델에 따르면, 의사결정자는 보상이 크지만 성공 확률이 희박한 게임에 참가하는 것을 거부할 수 있다. 특히, 사람들이 St. Petersburg 게임을 거부하는 이유는 참가자가 최종적으로 이익을 얻을 확률 $P_n(X, C)$가 0.5보다 낮기 때문이라고 해석하였다.

St. Petersburg 역설을 고려할 때, 일반적인 사람은 반복 시행 횟수 $n$을 무한대로 가정하지 않는다는 점을 알 수 있다. 예를 들어, 의사결정자가 고려하는 반복 시행 횟수를 $n=10,000$으로 설정할 경우, $X^{(2)}$ 게임에 비용 9를 지불하고 참가하면 이익 확률이 1에 가까워 합리적 선택으로 간주된다. 마찬가지로 St. Petersburg 게임에 비용 10을 지불하고 참가하는 경우에도 이익 확률이 1에 가까워 합리적 선택으로 판단된다. 반면, $X^{(20)}$ 게임에 비용 9를 지불하거나 St. Petersburg 게임에 비용 32를 지불하고 참가하는 것은 이익 확률이 0에 가까워 비합리적 선택으로 간주된다.

실증적인 연구 없이 대부분의 사람은 $X^{(2)}$ 게임에 비용 9를 지불하고 참가할 의사가 있다고 가정하자. 이 사람들은 St. Petersburg 게임에서 참가 비용이 10이라면 이익이 될 것이라고 생각할 것이다. 이는 두 상황에서 이익 확률의 수렴 속도가 비슷하기 때문으로, 의사결정자가 상황에 따라 자신의 $n$을 변경하지 않는다고 가정하면 의사결정자는 두 상황을 동일한 기준으로 판단하게 된다. 더 나아가, 사람들이 실제로 고려하는 반복 시행 횟수 $n$을 파악할 수 있다면, St. Petersburg 게임에서 이익 확률이 0.5가 되는 $C$를 계산함으로써 대부분의 사람이 합리적이라고 판단하는 적정 참가 비용을 도출할 수 있을 것이다.

\textbf{모델의 일반화} 시행 횟수-이익 확률 모델 $\mathcal{M}_n(X, C)$는 이익 확률을 최대화하는 전략을 확장하여, 의사결정자의 고유한 파라미터 $n$을 추가한 모델로 이해할 수 있다. 이 모델을 더욱 일반화하면, 의사결정자를 단일 값 $n$으로 표현하는 대신, $n$에 따른 가중치 곡선 $W_n$으로 나타내는 모델을 만들 수 있다. 이러한 접근에서는 $\sum_{n=1}^{\infty} P_n(X, C) W_n \overset{?}{=}0.5$ 를 의사결정 모델으로 설정할 수 있다.

\section*{모델의 한계}

첫째, 시행 횟수-이익 확률 모델 $\mathcal{M}_n(X, C)$는 개인화된 모델로, 특정 개인의 주관적인 판단 이유를 설명할 수 있다. 그러나 이 모델만으로는 St. Petersburg 게임의 객관적이고 공정한 참가 비용을 도출할 수 없다.

둘째, Lopes의 지적처럼 이 모델만으로는 객관적이고 합리적인 의사결정을 내리는 것이 불가능하다. 대신 이 모델을 의사결정의 보조 도구로 이용한다면, 기대효용이론을 적용해 장기적으로 손익을 평가한 뒤 개인화된 $\mathcal{M}_n(X, C)$를 사용하여 단기적으로 이익을 얻을 확률을 고려하는 방식을 생각해볼 수 있다. 또한 Samuelson의 지적대로 $\mathcal{M}_n(X, C)$는 이익 확률을 고려하는 모델이기 때문에 선택 간 전이성이 성립하지 않을 것이다. 모델이 제시한 선택 간 전이성이 성립하지 않는다면, 모델은 최적의 선택을 체계적으로 결정하지 못하는 한계가 있다.

\section*{논문의 한계}

본 논문에서는 $P_n(X, C)$의 수렴성을 중심극한정리를 활용하여 증명하였으나, 기댓값과 분산이 무한한 상황에 대해서는 수렴성을 증명하지 못하였다.

의사결정자가 고려하는 반복 시행 횟수 $n$이 $\mathcal{M}_n(X, C)$의 핵심적인 요소임에도 불구하고, 논문에서는 $n$의 본질적 의미와 개인이 $n$을 가지게 되는 과정을 설명하지 못하였다. 이 모델이 결정 이론의 연구들에도 부합하는 모델이라면, $n$을 개인의 위험 회피 성향과 시간 및 재산의 제약을 바탕으로 설명할 수 있을 것이다.

논문은 $\mathcal{M}_n(X, C)$를 활용하여 St. Petersburg 역설에 대한 해결책을 제시하는 데 초점을 맞췄으나, 해당 모델이 보다 일반적인 문제들에 대해 적용 가능함을 보이지 못하였다. 특히 $\mathcal{M}_n(X, C)$는 주어진 게임과 참가 비용에 대해 참가 여부만을 결정하는 단순화된 모델이다. 이 모델이 현실에서 의미를 가지려면 선택 간 전이성을 증명하고, 여러 선택지 중 최적의 선택을 도출할 수 있는 방법론을 제시해야 한다.

논문에서 다룬 게임들은 모두 $n$이 작을 때 이익 확률이 0에 가까웠다가, $n$이 증가함에 따라 이익 확률이 수렴하는 경우에 해당한다. 그러나 반대로 $n$이 작을 때 이익 확률이 1에 가까웠다가 $n$이 증가하면서 이익 확률이 0으로 수렴하는 경우에는 현실에서 사람들이 어떤 의사결정을 할 지 연구가 부족하다.
